% 三视图:球置于长方体之上;正:圆直径3(半径1.5)与矩形3×2相切,左侧直径3与高2及底宽3标注;侧:同形无标注;俯:正方形3×3内切圆;全实线
\begin{tikzpicture}[line width=.6pt,font=\normalsize,x=1cm,y=1cm]
% 正视图
\begin{scope}[shift={(0,4.8)}]
\draw (0,0) rectangle (3,2);
\draw (1.5,3.5) circle (1.5);
% 左侧标注:圆直径3(2~5),矩形高2(0~2)
\draw[thin] (0,0) -- (-0.5,0);
\draw[thin] (0,2) -- (-0.5,2);
\draw[thin] (0,2) -- (-0.5,2);
\draw[thin] (0,5) -- (-0.5,5);
\draw[<->] (-0.35,0) -- node[left,fill=white,inner sep=1pt] {$2$} (-0.35,2);
\draw[<->] (-0.35,2) -- node[left,fill=white,inner sep=1pt] {$3$} (-0.35,5);
% 底宽标注 3
\draw[thin] (0,0) -- (0,-0.35);
\draw[thin] (3,0) -- (3,-0.35);
\draw[<->] (0,-0.2) -- node[below,fill=white,inner sep=1pt] {$3$} (3,-0.2);
\node[below] at (1.5,-0.6) {正视图};
\end{scope}
% 侧视图
\begin{scope}[shift={(4.6,4.8)}]
\draw (0,0) rectangle (3,2);
\draw (1.5,3.5) circle (1.5);
\node[below] at (1.5,-0.6) {侧视图};
\end{scope}
% 俯视图(右,与正侧同高横向并排):正方形3×3,内切圆半径1.5
\begin{scope}[shift={(9.2,4.8)}]
\draw (0,0) rectangle (3,3);
\draw (1.5,1.5) circle (1.5);
\node[below] at (1.5,-0.6) {俯视图};
\end{scope}
\end{tikzpicture}
